\begin{mistake}{2026-06-29}{02}

\begin{problemcontent}
计算不定积分：
\[ \int \frac{x+5}{x^2 - 6x + 13} \mathrm{d}x \]
\end{problemcontent}

\begin{solutioncontent}
分母求导得 \( (x^2 - 6x + 13)' = 2x - 6 \)。
将分子凑出该形式：
\[ x+5 = \frac{1}{2}(2x - 6) + 8 \]
拆分积分：
\[
\begin{aligned}
\int \frac{x+5}{x^2 - 6x + 13} \mathrm{d}x &= \int \frac{\frac{1}{2}(2x - 6) + 8}{x^2 - 6x + 13} \mathrm{d}x \\
&= \frac{1}{2} \int \frac{2x - 6}{x^2 - 6x + 13} \mathrm{d}x + \int \frac{8}{x^2 - 6x + 13} \mathrm{d}x \\
&= \frac{1}{2} \int \frac{\mathrm{d}(x^2 - 6x + 13)}{x^2 - 6x + 13} + \int \frac{8}{(x-3)^2 + 2^2} \mathrm{d}x \\
&= \frac{1}{2} \ln(x^2 - 6x + 13) + 8 \cdot \frac{1}{2} \arctan \frac{x-3}{2} + C \\
&= \frac{1}{2} \ln(x^2 - 6x + 13) + 4 \arctan \frac{x-3}{2} + C
\end{aligned}
\]
（注：因为 \(x^2 - 6x + 13 > 0\) 恒成立，对数绝对值可直接写为圆括号）
\end{solutioncontent}

\mistakeknowledge{
\begin{itemize}
  \item \textbf{核心公式}：\(\int \frac{f'(x)}{f(x)} \mathrm{d}x = \ln |f(x)| + C\)，\(\int \frac{1}{x^2 + a^2} \mathrm{d}x = \frac{1}{a} \arctan \frac{x}{a} + C\)。
  \item \textbf{易错点}：拆分分子时，易漏掉常数项配方产生的反正切部分；或在计算反正切积分时漏乘系数 \(\frac{1}{a}\)（本题中为 \(\frac{1}{2}\)）；漏写积分常数 \(C\)。
  \item \textbf{题型信号}：分子为一次式，分母为无实根的二次多项式，标准解法必定为“先凑微分求对数，后配方求反正切”。
\end{itemize}}

\end{mistake}
